% ==============================================================================
% APA 7 Journal-Submission Manuscript Template
% ==============================================================================
% Copy to paper/main.tex. Standard: .claude/rules/working-paper-format.md
% Build:  cd paper && latexmk main.tex      (pdfLaTeX + biber via paper/latexmkrc)
%
% Class options:
%   man           professional manuscript (double-spaced, running head, title page)
%   floatsintext  tables/figures after first mention; REMOVE if the journal wants
%                 them after the references
%   mask          add for masked (blind) review -- hides authors and author note
% ==============================================================================
\documentclass[man,floatsintext]{apa7}

% ====== Language and Quotation (required by biblatex-apa) ======
\usepackage[american]{babel}
\usepackage{csquotes}

% ====== Encoding and Typography ======
\usepackage[T1]{fontenc}
\usepackage{microtype}

% ====== Math ======
\usepackage{amsmath, amssymb, mathtools}
\usepackage{bm}

% ====== Tables and Figures ======
\usepackage{booktabs}
\usepackage{threeparttable}
\usepackage{array, makecell, multirow}
\usepackage{siunitx}
\usepackage{graphicx}
\usepackage{pdflscape}
\graphicspath{{figures/}}

% ====== Lists ======
\usepackage{enumitem}

% ====== Bibliography (biblatex-apa + biber) ======
\usepackage[style=apa, sortcites=true, sorting=nyt, backend=biber]{biblatex}
\DeclareLanguageMapping{american}{american-apa}
\addbibresource{../Bibliography_base.bib}

% ====== Hyperref (second-to-last, options via \hypersetup) ======
\usepackage{hyperref}
\hypersetup{hidelinks, breaklinks=true}

% ====== Cleveref (immediately after hyperref) ======
\usepackage[capitalise, noabbrev, nameinlink]{cleveref}

% ==============================================================================
% Title page
% ==============================================================================
\title{[Paper Title in Title Case, About 12 Words or Fewer]}
\shorttitle{[RUNNING HEAD, 50 CHARACTERS OR FEWER]}

\authorsnames[1,2]{[Author One], [Author Two]}
\authorsaffiliations{
  {[Department of Psychology, University One]},
  {[School of Education, University Two]}}

\authornote{
  \addORCIDlink{[Author One]}{0000-0000-0000-0000}

  [Disclosures and acknowledgments: This study was preregistered at {registry
  URL}; deviations from the preregistration are reported in the Method section.
  Data, materials, and analysis code are available at {repository URL}. We have
  no known conflicts of interest to disclose. This research was supported by
  {funder, grant number}.]

  Correspondence concerning this article should be addressed to [Author One],
  [Department, University, Street Address, City, State Postal Code].
  Email: [email]}

\abstract{[One paragraph, 250 words or fewer (or the journal limit). State the
problem; the participants and their key characteristics (N, age or grade,
setting); the design and method; the main findings with effect sizes and 95\%
confidence intervals; and the conclusions and implications.]}

\keywords{[keyword one], [keyword two], [keyword three]}

% ==============================================================================
\begin{document}
\maketitle
% The class repeats the title here. The introduction has NO heading.

% \input{sections/introduction}
[Introduction: open with the problem, review the literature that motivates the
hypotheses, and close with the hypotheses or research questions. Example
citations: narrative \textcite{Simmons2011_false_positive}; parenthetical
\parencite{Appelbaum2018_jars_quant}.]

\section{Method}
\label{sec:method}
% \input{sections/method}

\subsection{Transparency and Openness}
[Preregistration (registry, date, link) and deviations; data, materials, and code
availability; reporting standard followed (JARS-Quant;
\cite{Appelbaum2018_jars_quant}); software and package versions.]

\subsection{Participants}
[Recruitment, eligibility, setting, dates of data collection, final sample with
demographic characteristics and how they were collected; exclusions and
attrition with reasons; ethics approval and consent/assent.]

\subsection{Sample Size Determination}
[A priori power analysis: target effect size and its source, alpha, power, and
for nested designs the ICC and cluster sizes (minimum detectable effect size).]

\subsection{Measures}
[Each measure: construct, source, number of items, response scale, scoring, and
reliability in this sample (coefficient omega, with CI).]

\subsection{Procedure}
[Assignment method and level of randomization, conditions, materials, timing,
fidelity monitoring.]

\subsection{Data Analysis}
[Models for each hypothesis, estimator, handling of nesting and missing data,
inference criteria, corrections for multiple comparisons; label exploratory
analyses.]

\section{Results}
\label{sec:results}
% \input{sections/results}

[Preliminary analyses, then confirmatory tests in preregistered order. Report
statistics per INV-4, e.g., \textit{t}(118) = 2.45, \textit{p} = .016,
\textit{d} = 0.45, 95\% CI [0.08, 0.81]. Refer to \Cref{tab:descriptives} and
\Cref{fig:example}.]

\begin{table}[tbp]
\begin{threeparttable}
\caption{Means, Standard Deviations, and Correlations Among Study Variables}
\label{tab:descriptives}
% \input{tables/descriptive/sumstats_correlations.tex}
\begin{tabular}{lcccc}
\toprule
Variable & \textit{M} & \textit{SD} & 1 & 2 \\
\midrule
1. [Variable one] & 0.00 & 1.00 & (.90) &       \\
2. [Variable two] & 0.00 & 1.00 & .00   & (.90) \\
\bottomrule
\end{tabular}
\begin{tablenotes}[para, flushleft]
{\small
\textit{Note.} $N = [n]$. Coefficient omega reliabilities appear in parentheses on
the diagonal. [Data source; missing-data handling.]
}
\end{tablenotes}
\end{threeparttable}
\end{table}

\begin{figure}[tbp]
\caption{[Figure Title in Title Case]}
\label{fig:example}
% \includegraphics[width=\linewidth]{fig1_description.pdf}
\centering
\fbox{\parbox[c][0.3\linewidth][c]{0.9\linewidth}{\centering [Figure placeholder]}}
\figurenote{[What is plotted, what error bars or bands represent (e.g., 95\% CIs),
and the data source.]}
\end{figure}

\section{Discussion}
\label{sec:discussion}
% \input{sections/discussion}

[Restate support for each hypothesis with effect sizes; interpret relative to
prior work; limitations; constraints on generality (populations, settings,
measures, and times to which the findings should and should not generalize);
implications for theory, practice, and policy.]

\printbibliography

% ====== Appendices (optional) ======
% \appendix
% \section{[Appendix Title]}
% \label{app:materials}

\end{document}
